\documentclass{article}
\usepackage[UTF8]{ctex}
\usepackage{amsmath,amssymb,amsfonts}
\usepackage{enumitem}
\usepackage{geometry}
\geometry{a4paper, margin=1in}
% === 错题本排版优化（不参与题目标识） ===
\setlength{\parindent}{0pt}
\setlength{\parskip}{0.35em plus 0.10em minus 0.08em}
\setlength{\emergencystretch}{3em}
\setlist{itemsep=0.25em, topsep=0.40em, parsep=0pt}
\allowdisplaybreaks[2]
% === 错题本排版优化结束 ===
\begin{document}
\section*{北京理工大学2025年数学分析真题}
\section*{4．北京理工大学2025年数学分析试卷（回忆版）}

\subsection{题目1}
1． \(\lim _{t \rightarrow 0} \frac{f(t)-2}{t}=3\) ，求 \(\lim _{t \rightarrow 0} \frac{e^{f(t)}-2}{t-2}\) ．

\subsection{题目2}
2． \(\lim _{x \rightarrow 0}\left(\frac{1}{\sin ^{2} x}-\frac{1}{x^{2}}\right)\) ．

\subsection{题目3}
3．\(f(x, y)=7 x y\) 在 \((x-1)^{2}+y^{2}=1\) 上的最大值．

\subsection{题目4}
4． \(\iint_{D}(x+2 x y) \mathrm{d} x \mathrm{~d} y\) ，其中 \(D:-1 \leq x \leq 1,-1 \leq y \leq x\) ．

\subsection{题目5}
5．\(F(x)=\int_{1}^{x^{2}} \mathrm{~d} t \int_{t^{2}}^{x^{2}} f(s, t) \mathrm{d} s\) ，求 \(F^{\prime}(x), F^{\prime \prime}(x)\) ．

\subsection{题目6}
6．\(u=x^{2} y+y^{2} z+z^{2} x\) ，求 \(\left.\frac{\partial^{2} u}{\partial x^{2}}\right|_{(0,0,1)},\left.\frac{\partial^{2} u}{\partial x \partial y}\right|_{(1,0,-1)},\left.\frac{\partial^{2} u}{\partial x \partial z}\right|_{(1,-1,-1)}\) ．

\subsection{题目7}
1．证明：\(e^{-x}+e^{x}>x^{2}-2, x \neq 0\) ．

\subsection{题目8}
2．\(f(x)\) 连续， \(\lim _{x \rightarrow 0} \frac{f(5 x-x)}{x}=3\) ，证明 \(f\) 在 \(x=0\) 处可导，且 \(f^{\prime}(0)=\frac{3}{4}\) ．

三．求幂级数 \(\sum_{n=0}^{\infty} \frac{x^{2 n+1}}{2 n+1}\) 的收敛半径，收敛点集，和函数．

四．\(u=f\left(\ln \sqrt{x^{2}+y^{2}}\right), \frac{\partial^{2} u}{\partial x^{2}}+\frac{\partial^{2} u}{\partial y^{2}}=\left(x^{2}+y^{2}\right)^{\frac{3}{2}}\) ，求 \(f(t)\) 的表达式．

五．\(f(x)=\sin \frac{2}{x}\) ，证明 \(f(x)\) 在 \((a,+\infty)(a>0)\) 上一致连续，在 \((0,+\infty)\) 上不一致连续．

六．\(f(x)\) 非负二阶可导，且存在 \(M>0\) 满足 \(\left|f^{\prime \prime}(x)\right| \leq M\) ，证明：\(f^{\prime}(x) \leq \sqrt{2 M f(x)}\) ．

七．已知 \(\lim _{x \rightarrow \infty} f(x)=f(0)\) ．

\begin{itemize}
  \item[（1）] 证明：存在 \(\xi_{1}, \xi_{2} \in(0,+\infty), \xi_{1} \neq \xi_{2}\) ，使得 \(f\left(\xi_{1}\right)=f\left(\xi_{2}\right)\) ．
  \item[（2）] \(f(x)\) 至少有一个最大值或最小值．
  \item[（3）] \(0 \leq f(x) \leq \frac{x^{3}}{e^{x}}\) ，证明：存在 \(\xi\) 满足 \(f^{\prime}(\xi)=\frac{\xi^{2}(3-\xi)}{e^{\xi}}\) ．
\end{itemize}

八．\(I(p)=\int_{0}^{+\infty} \frac{\cos ^{2} x}{x^{p}} \mathrm{~d} x\) ．

\begin{itemize}
  \item[（1）] 求 \(I(p)\) 的定义域．
  \item[（2）] 讨论 \(I(p)\) 在定义域上的连续性．
  \item[] \((3)\) 证明 \(I(p)\) 在 \((-1,1)\) 不一致收敛．
\end{itemize}
\end{document}
